\documentclass[11pt,a4paper]{article}

\usepackage[T1]{fontenc}
\usepackage[utf8]{inputenc}
\usepackage[italian]{babel}
\usepackage{lmodern}
\usepackage[margin=2.8cm]{geometry}
\usepackage{amsmath}
\usepackage{amssymb}
\usepackage{booktabs}
\usepackage{longtable}
\usepackage{microtype}
\usepackage{enumitem}
\usepackage[hidelinks]{hyperref}

\setlist[itemize]{itemsep=2pt,topsep=4pt}
\setlist[enumerate]{itemsep=2pt,topsep=4pt}

\title{Post di Trump e prezzo del greggio\\[2pt]
\large Stato del progetto}
\author{}
\date{21 settembre 2026}

\begin{document}
\maketitle
\vspace{-2em}

\section{La catena di elaborazione}

La domanda è se i post di Trump che segnalano un rischio all'offerta di
greggio muovano il prezzo del petrolio, e per quanto tempo. Rispondere
richiede tre cose: sapere \emph{quando} ogni post è stato pubblicato, con
precisione sufficiente a isolarne l'effetto; stabilire \emph{quali} post
contano come eventi, con un criterio che un revisore possa replicare; e avere
una serie di prezzo abbastanza fitta da misurare la reazione nei minuti
successivi.

Il lavoro, fin ora, è diviso in cinque script numerati, ciascuno dei quali legge
l'output del precedente.

\begin{description}[leftmargin=1.6cm,style=nextline,itemsep=5pt]

\item[\texttt{00}\ \ scraping] Costruisce il corpus dall'archivio
dell'American Presidency Project. Il punto non banale è il tempo: l'orario di
pubblicazione non viene letto dalla pagina ma ricavato dall'identificatore del
post, che lo contiene per costruzione. Il confronto fra i due serve da
controllo di qualità sull'intero dataset.

\item[\texttt{01}\ \ selezione] Porta da ``tutti i post'' a ``gli eventi
rilevanti''. Scarta ciò che non ha contenuto testuale, applica la regola che
definisce l'evento petrolifero, raggruppa i post ravvicinati in episodi e
misura quanto possono essere lunghe le finestre di osservazione prima che due
episodi si sovrappongano. L'episodio, non il post, è l'unità di analisi.

\item[\texttt{02}\ \ audit] Risponde all'obiezione prevedibile a qualunque
definizione basata su parole chiave: \emph{hai trovato solo quello che
cercavi}. Applica due metodi non supervisionati, che non conoscono le liste di
termini, e verifica se emergono post rilevanti fuori dalla regola. Non
definisce nulla: serve a misurare quanto la regola perde.

\item[\texttt{03}\ \ prezzi] Procura la serie di prezzo al minuto contro cui
misurare la reazione. È la fase ferma, in attesa di decidere insieme quali
serie storiche usare.

\item[\texttt{04}\ \ feature] Porta da ``quali post contano'' a ``quanto e in
che direzione ciascuno spinge''. Feature calcolate da regole (orario,
lunghezza, enfasi), novità semantica rispetto agli episodi recenti, e le
variabili che richiedono lettura: direzione attesa sull'offerta, intensità,
concretezza.

\end{description}

I primi tre script costruiscono la variabile esplicativa e ne verificano la
tenuta; il quarto la gradua; il terzo fornisce la variabile dipendente. La
separazione è deliberata: le feature vanno congelate prima di guardare i
prezzi, altrimenti tararle sui rendimenti significherebbe costruire il
risultato invece di misurarlo.

\section{Dove siamo}

La pipeline arriva fino alle feature deterministiche per episodio. Testo,
regola di selezione, episodi e audit non supervisionato sono stati eseguiti e
hanno prodotto output su disco; i prezzi sono ancora allo stadio di test; la
parte di annotazione (manuale e LLM) non è iniziata.

\begin{center}
\begin{tabular}{llr}
\toprule
Fase & Stato & Output \\
\midrule
00 scraping UCSB        & eseguito & 8\,619 post, 98,1\% verificati \\
01 pulizia e regola     & eseguito, in taratura & 4\,575 post, 362 eventi, 329 episodi \\
02 audit BERTopic       & eseguito & 38 cluster, 84 candidati falsi negativi \\
03 prezzi Dukascopy     & solo test su marzo 2025 & 10 file in \texttt{download/} \\
04 feature              & solo fase A & 329 episodi, 18 colonne \\
\bottomrule
\end{tabular}
\end{center}

\section{Il corpus e i timestamp}

Lo scraper ha interrogato tutti i 608 giorni fra il 1 dicembre 2024 e il 31
luglio 2026. Tre pagine risultano assenti sul server, cinque giorni sono
elencati in \texttt{giorni\_falliti.txt}, i restanti hanno prodotto
8\,619 post distribuiti su 597 giorni (mediana 10 post al giorno, massimo
177).

Il timestamp è derivato dall'identificatore del post. Truth Social eredita da
Mastodon gli ID snowflake, in cui i bit alti sono i millisecondi dall'epoch
Unix:
\[
  t_{\text{UTC}} = \frac{\texttt{post\_id} \gg 16}{1000},
  \qquad
  \texttt{seq} = \texttt{post\_id} \;\&\; \texttt{0xFFFF}.
\]
Il campo di sequenza ordina i post pubblicati nello stesso millisecondo, cosa
che serve dentro i burst. Il 98,0\% dei timestamp ha secondi diversi da zero,
come atteso: si può essere precisi anche al millisecondo.

\subsection*{Il formato dell'ora in pagina}

L'archivio non usa un formato solo. Su un campione di 766 blocchi di post, 457
riportano l'ora a 24 ore senza suffisso (\texttt{18:11}) e 305 a 12 ore con
meridiano (\texttt{1:17 PM}, \texttt{7:05 A.M.}). La prima versione del parser
cercava \verb|(\d{1,2}):(\d{2})| e ignorava il suffisso, quindi leggeva
\texttt{1:17 PM} come l'una di notte.

L'effetto sulla validazione era vistoso: 1\,908 righe segnalate come anomale,
di cui 1\,752 con uno scarto di esattamente 720 minuti — cioè dodici ore, la
firma di un meridiano perduto. La correzione cerca l'ora, poi ancora la
ricerca del meridiano subito dopo di essa, e richiede maiuscola e confine di
parola per non scambiare per meridiano l'inizio del testo che segue un orario
a 24 ore (\texttt{10:08 Amazing\dots}, \texttt{09:36 A MUST WATCH\dots}: due
casi presenti nel corpus).

Dopo la correzione le righe verificate passano da 6\,711 a \textbf{8\,455 su
8\,619, il 98,1\%}. Le 164 anomalie residue si scompongono così:

\begin{center}
\begin{tabular}{lrl}
\toprule
Scarto & Righe & Causa \\
\midrule
180 min & 105 & UCSB stampa in Eastern time dal 15 al 21 settembre 2025 \\
altro   & 36  & scarti irregolari, errori redazionali isolati \\
assente & 15  & ora non presente in pagina \\
720 min & 8   & pagina a 24 ore con ora errata (\texttt{12:27} per \texttt{00:27}) \\
\bottomrule
\end{tabular}
\end{center}

I 105 casi a tre ore cadono tutti in una singola settimana, in cui l'archivio
ha cambiato fuso senza dichiararlo; gli 8 residui a 720 minuti non sono un
formato che sfugge alla regex ma errori di trascrizione su pagine che usano il
formato a 24 ore. Nessuno dei due gruppi tocca i timestamp effettivi, che
vengono dagli ID: riguardano solo la colonna di controllo. La settimana in
Eastern time resta segnalata come anomala, in attesa di decidere se trattarla
come eccezione dichiarata.

\section{La pulizia}

Dei 8\,619 post scaricati ne sopravvivono alla pulizia \textbf{4\,575}, il
53,1\%. 
Segue una tabella dei post rimossi:
\begin{center}
\begin{tabular}{lrl}
\toprule
Criterio & Post rimossi & Cosa toglie \\
\midrule
Segnaposto senza contenuto & 1\,849 & repost e immagini senza didascalia \\
Doppi invii entro 2 minuti &   497  & testi identici ravvicinati \\
Boilerplate endorsement    &   730  & il template degli endorsement elettorali \\
Sotto le 5 parole          &   968  & slogan e frammenti \\
\midrule
Totale                     & 4\,044 & \\
\bottomrule
\end{tabular}
\end{center}

\subsection*{La causa: un post su tre non ha testo proprio}

Trump pubblica moltissimi video e immagini senza scrivere nulla. Nella pagina
d'archivio un post del genere compare come la parola \emph{Video} seguita
dall'indirizzo del file, oppure come il solo indirizzo. Gli indirizzi vengono
rimossi all'inizio della pulizia — sono rumore per qualunque ricerca
lessicale — e quello che resta è una stringa vuota o la sola parola
\texttt{Video}. In tutto sono \textbf{2\,749 post, il 31,9\% del corpus}:
1\,804 diventano vuoti e 945 si riducono a \texttt{Video}.

Sono post reali, non errori di scraping, ma non contengono niente su cui una
regola basata su parole possa lavorare. Escluderli è inevitabile. 
\subsection*{Il criterio dei doppi invii rimuove post reali}

Il caso da correggere è il secondo filtro. Serve a eliminare il doppio invio
accidentale: lo stesso messaggio pubblicato due volte a pochi secondi di
distanza è un gesto comunicativo solo, non due eventi. Per riconoscere ``lo
stesso messaggio'' i testi vengono confrontati in forma \emph{normalizzata}:
tutto minuscolo, via la punteggiatura, spazi compattati. È una chiave di
confronto, non il testo che verrà analizzato dopo. Così ``Deal with Iran!!!''
e ``deal with iran'' diventano entrambi \texttt{deal with iran} e risultano
uguali.

\subsection*{Cosa comporta}
Va sottolineato che il 47\% dei post non è scartato in fase di pulizia per una 
selezione aggressiva.
I post sono scartati perché privi di contenuto e quindi per "costruzione", 
sono irrilevanti per i nostri scopi.

Si può valutare l'implementazione di un filtro a parte più specifico che chiarisca
con precisione le proporzioni di tali post per un puro fine analitico. Ai fini 
dell'estrazione delle features è irrilevante.

Resta un limite che non si risolve spostando filtri. Un post di solo media può
benissimo essere rilevante per la domanda di ricerca — un video di una
raffineria colpita è informazione di mercato a tutti gli effetti — e nessuna
regola lessicale può riconoscerlo. Conviene dichiararlo insieme al criterio di
campionamento invece di lasciarlo implicito in una fase di pulizia.

\section{La regola e gli eventi}

La regola richiede la congiunzione di un attore dell'offerta e di un
meccanismo:
\[
  \text{evento}_i = \mathbb{1}\!\left[\,
    |A \cap T_i| \ge 1 \;\wedge\; |M \cap T_i| \ge 1 \,\right],
\]
con $T_i$ i token del post $i$, $A$ la lista degli attori e $M$ quella dei
meccanismi. Sul corpus pulito: 631 post con almeno un attore, 810 con almeno
un meccanismo, \textbf{362 eventi} (7,9\%).  I termini che
reggono la regola sono concentrati — \texttt{iran} (223), \texttt{russia}
(65), \texttt{hormuz} (62), \texttt{ukraine} (57) fra gli attori;
\texttt{military} (97), \texttt{war} (95), \texttt{nuclear} (87),
\texttt{strait} (67), \texttt{oil} (64) fra i meccanismi.

Due problemi di taratura emergono dalle liste stesse.

\textbf{Termini di mercato usati come attori.} \texttt{market} e
\texttt{markets} sono nella lista degli attori, e \texttt{boom},
\texttt{come down}, \texttt{going down}, \texttt{coming down} in quella dei
meccanismi. Non sono attori dell'offerta né meccanismi di interruzione: sono
commento sui prezzi. Venticinque eventi (6,9\%) si reggono esclusivamente su
questi termini.
Anche se leggendoli sono chiaramente fuori tema (e.g. ``The MARKETS are
going to BOOM'', ``This is Biden's Stock Market'') li ho inseriti perché le 
date di molti post sono post attacco in Venezuela o in iran da interpretare 
come un riferimento a quanto stava succedendo in quel periodo. 
Se tolti, gli eventi scendono a 337. Il costo in numerosità è modesto.

\textbf{Il 44,5\% degli eventi si regge solo su meccanismi deboli}
(\texttt{war}, \texttt{military}, \texttt{attack}, \texttt{strike},
\texttt{nuclear}, \texttt{down}), cioè termini che indicano conflitto ma non
offerta di greggio. La diagnostica è già implementata e segnala il caso; 161
eventi su 362 ricadono lì. 
Nel file \texttt{dati/petrolio/da\_leggere.csv} è presente un campione di 150 da controllare per verificare
la correttezza della regola. Da una prima esamina si capisce che la regola tende più 
a prendere un falso positivo che un falso negativo (Trump a volte tira fuori 
le parole chiavi scelte solo per parlare male dei suoi avversari politici).
Tutta via a causa della già scarsa presenza di post ho preferito non proseguire
con una cernita più severa. 

\textbf{UPDATE}-- \\
Ho creato la cartella annotazioni che tiene traccia delle varie versioni della regola e dei file da controllare. 
questa sarà utile sia a tenere traccia delle performance di ogni versione della regola sia ad autocompilare il giudizio dei post che sono già stati giudicati.

\subsection{Il criterio della regola}
si ritengono interessanti i post in cui Trump parla di: 
\begin{enumerate}
  \item \textbf{il conflitto iraniano e lo stretto di hormuz}: qualsiasi news riguardi l'apertura dello stretto o il raggiungimento di un accordo con gli iraniani è rilevante per capire le aspettative future sulla rotta petrolifera più importante. 
  \item \textbf{Venezuela}: l'attacco intrapreso dall'amministrazione Trump in Venezuela ha interessato gli approvvigionamenti di petrolio americano dalle navi venezuelane 
  \item \textbf{Previsioni di Trump sul mercato}: in molti post trump incita i suoi lettori a comprare. Spesso si stratta di propaganda e basta. Ciò nonostante, in un'analisi che riguarda l'influenza che ha un uomo sul mercato è interessante anche quando il presidente USA elargisce esplicitamente consigli finanziari. 
  
\end{enumerate}

\section{Audit non supervisionato}

L'audit è stato rilanciato sul corpus corretto: 38 cluster con il
37,6\% di outlier. La copertura della regola per cluster è coerente con il
disegno: 0,66 sul cluster Iran/Hormuz (195 su 297 post catturati), 0,65 su
Ucraina/Russia (46 su 71). Le piccole variazioni rispetto all'esecuzione
precedente vengono dalla ricostruzione degli embedding su testi ripuliti e non
cambiano il quadro.

Tre cluster hanno lessico petrolifero e copertura bassa. Il più rilevante è
Venezuela/\allowbreak oil/\allowbreak strike, con copertura \textbf{0,43}: 23
post catturati su 54. Il tema è dichiaratamente dentro il perimetro della
tesi, quindi quella copertura indica un buco vero nella lista dei termini, non
un falso allarme.

Gli altri due stanno fuori per costruzione: uno raccoglie la politica
energetica interna (EPA, permessi), l'altro — \texttt{prices},
\texttt{gasoline}, \texttt{inflation} — il prezzo al distributore, con
copertura \textbf{0,00}. Che il clustering isoli quest'ultimo a ogni
riesecuzione non è un difetto della regola, ma dice che il confine fra prezzo
al consumo americano e offerta globale di greggio conviene scriverlo per
esteso invece di lasciarlo implicito nelle liste.

La ricerca semantica restituisce 84 candidati falsi negativi, con similarità
mediana 0,45 e massima 0,65. I più vicini alle sonde sono commenti sul prezzo
al distributore (e.g. ``Gasoline Retailers must get their Prices down'',
``Gasoline just broke \$1.98 a Gallon''). Nessuna
parafrasi di interruzione dell'offerta compare fra i candidati: il richiamo
della regola sul suo perimetro dichiarato sembra alto, e questo è il risultato
che l'audit doveva produrre.

\section{Episodi, finestra evento, potenza}

Post entro 30 minuti l'uno dall'altro sono collassati in un episodio, datato
al primo post. I 362 eventi danno \textbf{329 episodi}: 297 con un solo post,
31 con due, 1 con tre. Solo il 9,7\% è multi-post, il che rende
l'attribuzione del movimento molto più semplice di quanto ci si aspettasse.
Identità e datazione degli episodi sono invariate rispetto a prima della
correzione: cambia il testo di 179 di essi, non quali siano.

32,5\% in orario di mercato USA, 24,0\% nel weekend, 16,9 episodi al mese su un arco di
593 giorni. L'addensamento è marcato — giugno 2025, marzo, aprile e giugno
2026 — e fornisce i regimi naturali da confrontare con CD-NOD senza doverli
scegliere a posteriori.

Gli intervalli fra episodi sono ampi: il primo percentile è 36 minuti, la
mediana 19,5 ore. La quota di episodi contaminati da un altro dentro la
finestra resta nulla fino a 30 minuti, sale al 3,4\% a un'ora e al 10,1\% a
due ore. Con la soglia del 15\% adottata, \textbf{la finestra evento arriva a
due ore} senza problemi di sovrapposizione. Il vincolo, quindi, non viene dal
testo ma dai prezzi.

La numerosità minima rilevabile si legge dalla formula standard per un test a
due code su una media:
\[
  \text{MDE} = \frac{z_{1-\alpha/2} + z_{1-\beta}}{\sqrt{n}}\,\sigma
  \;=\; \frac{2{,}802}{\sqrt{n}}\,\sigma
  \qquad (\alpha = 0{,}05,\ 1-\beta = 0{,}8).
\]
Con $n = 329$ si rilevano effetti da 0,154 deviazioni standard del rendimento
sulla finestra. Dividendo per direzione attesa si sale attorno a 0,22, e sui
247 episodi che avranno effettivamente un prezzo (sezione seguente) a 0,178.
Sono soglie praticabili per un evento macro su finestre di minuti, ma non
lasciano margine per sottogruppi ulteriori.

\section{Dati di prezzo: cosa dicono i test}

\begin{center}
\begin{tabular}{lrrr}
\toprule
& Brent (\texttt{brentcmdusd}) & WTI (\texttt{lightcmdusd}) & Oro (\texttt{xauusd}) \\
\midrule
Barre nel mese          & 12\,720 & 36\,780 & --- \\
Giorni feriali coperti  & 10 / 21 & 21 / 21 & 5 / 5 nel test \\
Ore UTC coperte         & 0--20   & 0--23   & 0--23 \\
Barre per giorno        & 1\,260  & 1\,440  & 1\,440 \\
\bottomrule
\end{tabular}
\end{center}

Il Brent è inutilizzabile: metà dei giorni feriali manca del tutto e i giorni
presenti si fermano alle 20:59 UTC. Il WTI copre tutti i giorni feriali e
l'intera giornata. Ristretto a lunedì--venerdì: nessun buco superiore a 15
minuti, 5,5\% di barre a volume nullo, deviazione standard dei rendimenti al
minuto di $4{,}3 \times 10^{-4}$, 20\% di rendimenti esattamente nulli. Sono
valori normali per dati al minuto su un future energetico.

La riserva riguarda la domenica. Le barre domenicali sono il 92,6\% a volume
nullo e riportano il prezzo di chiusura del venerdì congelato per ventidue
ore: cinque tratti da oltre 1\,300 minuti a prezzo identico, uno per ogni
domenica del mese. Sono riempimento sintetico, non quotazioni, e vanno
scartate. Sul totale del mese il 21\% delle
barre cade in tratti costanti di almeno un'ora; sui soli giorni feriali
scende al 3,6\%.

Tenendo conto di sabato, domenica fino alle 22 UTC e venerdì dopo le 21,
\textbf{82 episodi su 329 (24,9\%) restano senza quotazione} e ne rimangono
247 utilizzabili. Con il solo Brent sarebbero stati 211. La perdita non è
casuale: colpisce un'intera fascia di calendario, quindi va confrontata la
distribuzione delle feature fra episodi persi e conservati. Il primo controllo
disponibile — il numero di post per episodio — non mostra differenze (1,08
contro 1,10), ma il confronto va rifatto su tutte le feature una volta che
esistono.

Da qui la fase è ferma. Il test dice che una strada praticabile esiste, non
che sia quella giusta, e la scelta conviene farla prima di scaricare venti
mesi di dati. Restano aperte tre questioni.

\begin{enumerate}
  \item \textbf{Quanto in là arriva il WTI.} Marzo 2025 è coperto senza buchi,
    ma è un mese solo. La copertura sul resto del periodo non è verificata, e
    lo stesso vale per l'oro, scaricato finora su cinque giorni.
  \item \textbf{Che dato è.} Sono quotazioni CFD di un broker, non prezzi di
    borsa: seguono il future sottostante con uno spread proprio. Su finestre
    di pochi minuti la differenza dovrebbe essere trascurabile, ma è
    un'assunzione da avallare, non da dare per buona.
  \item \textbf{Gli episodi senza mercato.} Un quarto degli episodi cade
    quando non esiste un prezzo. Si accetta la perdita dichiarandola, si cerca
    uno strumento con copertura più ampia, o si tratta quella fascia in modo
    diverso nel disegno?
\end{enumerate}

Lo script 03 è ancora impostato sul Brent come variabile principale: verrà
allineato quando la scelta sarà fatta.

\section{Feature}

\texttt{episodi\_features.parquet} contiene i 329 episodi con 18 colonne.
Sono presenti solo le feature deterministiche della fase A — ora UTC e New York, giorno, weekend,
orario di mercato, caratteri, esclamativi, maiuscole, parole urlate.

La correzione del parser ha avuto un effetto misurabile. Le due
lettere \texttt{AM}/\texttt{PM} contavano come maiuscole e gonfiavano la
misura di enfasi: la quota di maiuscole era sovrastimata su \textbf{177
episodi su 329}, di 0,005 in mediana e fino a 0,027. La mediana complessiva
scende da 0,1065 a 0,1029. È una distorsione piccola ma non innocua, perché
dipendeva dal formato della pagina d'archivio e non dal post, quindi entrava
nelle feature come rumore arbitrario.

Mancano la novità semantica, il gold standard e lo scoring LLM. La novità è
l'unica delle tre che non richiede annotazione:
\[
  \text{novit\`a}_i = 1 - \max_{\,j\,:\,t_i - 30\text{gg} \le t_j < t_i}
  \cos\!\left(e_i, e_j\right),
\]
con $e_i$ l'embedding del testo unito dell'episodio. Conviene calcolarla
subito, perché è la feature che serve a distinguere l'informazione nuova dalla
ripetizione ed è indipendente da tutto il resto.

Il codice della fase D salva già modello, prompt e temperatura accanto ai
risultati, che è la condizione minima perché una feature generata da un LLM
sia riproducibile.

\section{Decisioni da prendere insieme}

Prima di proseguire ho bisogno di un consiglio sulle seguenti tematiche: 

\begin{enumerate}
  \item \textbf{La serie di prezzo}, nei termini della sezione precedente.
    Finché non è decisa, la variabile dipendente non esiste e l'event study
    non si può impostare.
  \item \textbf{La settimana in Eastern time.} Dal 15 al 21 settembre 2025
    l'archivio stampa gli orari in un fuso diverso senza dichiararlo. I
    timestamp restano corretti perché vengono dagli identificatori, ma il
    controllo di qualità segnala 105 righe. O diventa un'eccezione dichiarata
    nel validatore, o resta un'anomalia documentata: in tesi cambia solo come
    la si racconta.
  \item \textbf{Il confine del tema.} Un post sul prezzo della benzina al
    distributore non è un rischio all'offerta globale, ma è a un passo, e sia
    la ricerca semantica sia il clustering continuano a riportarlo a galla.
    Conviene fissare il criterio adesso e per iscritto, prima che sia
    un'annotazione a deciderlo caso per caso.
\end{enumerate}

\section{Lavoro già pianificato}

\begin{enumerate}
  \item Introdurre un filtro esplicito per i post di solo media, contato a
    parte, invece di lasciarli distribuiti su tre criteri; e rendere il
    criterio dei doppi invii condizionale alla presenza di testo.
  \item Valutare di togliere \texttt{market}, \texttt{markets} e i termini di prezzo dalle
    liste; aggiungere i termini mancanti del cluster Venezuela.
  \item Calcolare la novità semantica ed esportare il gold standard;
    annotarlo, eseguire lo scoring LLM, misurare il kappa pesato.
  \item Congelare le feature e annotare la data.
\end{enumerate}

\end{document}
